\metatitle{Varieties in algebraic combinatorics}
\metadescription{A short guide to Grassmannians, flag varieties, Schubert varieties, and related spaces.}

\section[varieties]{Varieties in algebraic combinatorics}

Many constructions in algebraic combinatorics are shadows of geometric
objects.  A polynomial may be the character of a representation, the class of
a subvariety, the Hilbert series of a coordinate ring, or the point-counting
shadow of a variety over a finite field.  This page gives a short map of the
varieties that occur most often on this site.

The main examples are \hyperref[grassmannianVarieties]{Grassmannians},
\hyperref[flagVarieties]{flag varieties},
\hyperref[schubertVarietiesAndMatroids]{Schubert varieties},
\hyperref[richardsonVarieties]{Richardson varieties},
\hyperref[springerFibers]{Springer fibers}, and
\hyperref[hessenbergVarietiesBackground]{Hessenberg varieties}.  They connect
directly to \hyperref[schubertCalculus]{Schubert calculus},
\hyperref[grothendieck]{Grothendieck polynomials},
\hyperref[key]{key polynomials}, \hyperref[schurFlagged]{flagged Schur
polynomials}, \hyperref[positroid]{positroids}, and
\hyperref[matroid]{matroids}.


\section[affineProjectiveVarieties]{Affine and projective varieties}

An \defin{affine variety}\label{affineVarietyDefinition} is, roughly, the set
of common zeros of polynomial equations in affine space:
\[
  X=\{x\in\setC^n:f_1(x)=\dotsb=f_r(x)=0\}.
\]
A \defin{projective variety}\label{projectiveVarietyDefinition} is the
corresponding notion in projective space, where points are lines through the
origin:
\[
  \mathbb P^{n-1}
  =
  (\setC^n\setminus\{0\})/\setC^\ast.
\]
Projective varieties are the natural home for many compact geometric objects,
including Grassmannians and flag varieties.

The coordinate ring viewpoint is often combinatorial.  If $X$ is affine with
coordinate ring $\setC[X]$, then a grading on $\setC[X]$ gives a Hilbert
series.  Such Hilbert series often become symmetric functions, lattice-point
generating functions, or Ehrhart-type series after a suitable choice of
coordinates.


\section[grassmannianVarieties]{Grassmannians}

The \defin{Grassmannian}\label{grassmannianVarietyDefinition}
$\operatorname{Gr}(k,n)$ is the variety of $k$-dimensional subspaces of
$\setC^n$.  Equivalently, it parametrizes projective $(k-1)$-planes in
$\mathbb P^{n-1}$.  For example, $\operatorname{Gr}(2,4)$ parametrizes lines
in $\mathbb P^3$.

\begin{example}
The Grassmannian $\operatorname{Gr}(1,n)$ is projective space
$\mathbb P^{n-1}$, since a one-dimensional subspace of $\setC^n$ is exactly
a line through the origin.
\end{example}

The standard projective realization is the
\defin{Plücker embedding}\label{pluckerEmbedding}
\[
  \operatorname{Gr}(k,n)\hookrightarrow \mathbb P^{\binom nk-1}.
\]
If a $k$-plane is represented by a full-rank $k\times n$ matrix, its maximal
minors are the \defin{Plücker coordinates}\label{pluckerCoordinate}.  These
coordinates satisfy quadratic Plücker relations.

This is where matroids enter.  Given a point of $\operatorname{Gr}(k,n)$, the
sets $I\subseteq[n]$ with nonzero Plücker coordinate $p_I$ form the
\hyperref[matroidBasis]{bases} of a rank $k$ representable
\hyperref[matroidDefinition]{matroid}.  The totally nonnegative part of the
Grassmannian gives \hyperref[positroid]{positroids}, introduced by
\name{Alexander Postnikov} \cite{Postnikov2006x}.  Schubert matroids and
\hyperref[latticePathMatroidDefinition]{lattice path matroids} are important
positroid families \cite{Oh2011}.


\section[flagVarieties]{Flag varieties}

A \defin{complete flag}\label{completeFlagDefinition} in $\setC^n$ is a chain
\[
  0=F_0\subset F_1\subset F_2\subset\dotsb\subset F_n=\setC^n,
  \qquad \dim F_i=i.
\]
The \defin{flag variety}\label{flagVarietyDefinition} $\operatorname{Fl}_n$ is
the variety of all complete flags.  A
\defin{partial flag variety}\label{partialFlagVarietyDefinition} remembers only
some of the subspaces in the chain.  Grassmannians are partial flag varieties,
since
$\operatorname{Gr}(k,n)$ remembers only the $k$-dimensional subspace.

The flag variety is the geometric source of many divided-difference
constructions.  Its Schubert classes are indexed by permutations, and the
corresponding polynomial representatives are
\hyperref[schubert]{Schubert polynomials}.  Passing from cohomology to
\hyperref[kTheory]{K-theory} replaces these by
\hyperref[grothendieck]{Grothendieck polynomials}.


\section[schubertVarietiesAndMatroids]{Schubert varieties and Schubert matroids}

Fix a complete flag $F_\bullet$ in $\setC^n$.
\defin{Schubert varieties}\label{schubertVarietyGeometryDefinition} are
subvarieties of a Grassmannian or flag variety defined by incidence conditions
with respect to $F_\bullet$.  In the Grassmannian, they are indexed by
partitions in a rectangle, or equivalently by $k$-element subsets of $[n]$.

For $I=\{i_1<\dotsb<i_k\}\subseteq[n]$, the corresponding
\defin{Schubert matroid}\label{schubertMatroidDefinition} has bases
\[
  \mathcal B_I
  =
  \{J=\{j_1<\dotsb<j_k\}: i_a\leq j_a\text{ for all }a\}.
\]
For example, if $I=\{1,3\}\subseteq[4]$, then
\[
  \mathcal B_I=\{13,14,23,24,34\}.
\]
Thus Schubert matroids are the matroid-theoretic shadow of Schubert cells in
the Grassmannian.  They are positroids, and Suho Oh showed that positroids
can be characterized using Schubert matroids \cite{Oh2011}.

This viewpoint helps explain why matroid polytopes, positroid polytopes,
Schubert calculus, and Grassmannian coordinates often appear together.


\section[richardsonVarieties]{Richardson varieties}

A \defin{Richardson variety}\label{richardsonVarietyDefinition} is an
intersection of a Schubert variety and an opposite Schubert variety.  In the
complete flag variety, if $X_w$ is the Schubert variety indexed by $w$ and
$X^v$ is the opposite Schubert variety indexed by $v$, then
\[
  X_w^v\coloneqq X_w\cap X^v.
\]
This intersection is nonempty precisely when $v\leq w$ in
\hyperref[bruhatOrder]{Bruhat order}.

Richardson varieties are useful because they isolate the geometry of an
interval in \hyperref[bruhatOrder]{Bruhat order}.  They also appear naturally in total positivity,
positroid geometry, and the study of components of Springer fibers.  For a
recent tableau connection, see the Richardson tableaux of
\name{Steven N. Karp} and \name{Martha E. Precup} \cite{KarpPrecup2025x}.


\section[hessenbergVarietiesBackground]{Hessenberg varieties}

\defin{Hessenberg varieties}\label{hessenbergVarietyDefinition} are
subvarieties of the flag variety cut out by a linear operator and a
\defin{Hessenberg function}\label{hessenbergFunctionDefinition}.  In type $A$,
one starts with a weakly increasing sequence $m=(m_1,\dotsc,m_n)$ with
$i\leq m_i\leq n$ and considers flags satisfying
\[
  T F_i\subseteq F_{m_i}.
\]
The precise behavior depends on the choice of the operator $T$.

They were introduced by \name{F. De Mari}, \name{C. Procesi}, and
\name{M. A. Shayman} \cite{MariProcesiShayman1992}.  Regular semisimple
Hessenberg varieties are the varieties behind the
\hyperref[dotAction]{dot action} and the Shareshian--Wachs connection with
\hyperref[chromaticQuasisymmetricHessenberg]{chromatic quasisymmetric
functions}.


\section[springerFibers]{Springer fibers and the Peterson variety}

A \defin{Springer fiber}\label{springerFiberDefinition} is a subvariety of the
flag variety attached to a nilpotent linear operator $N$.  In type $A$, it is
\[
  \mathcal B_N
  =
  \{F_\bullet\in\operatorname{Fl}_n:N F_i\subseteq F_i
  \text{ for all }i\}.
\]
The cohomology of Springer fibers carries deep Weyl-group representations;
this is the starting point of Springer theory \cite{Springer1978}.  Springer
fibers also interact with tableaux, Hall--Littlewood theory, and the
\hyperref[deltaSpringerHallLittlewood]{$\Delta$-Springer varieties} appearing
near the Delta conjecture.

The \defin{Peterson variety}\label{petersonVarietyDefinition} is a special
regular nilpotent Hessenberg variety.  In type $A$, it is obtained from a
regular nilpotent operator $N$ and the Hessenberg function
\[
  m(i)=i+1\quad(1\leq i<n),\qquad m(n)=n.
\]
Thus its flags satisfy
\[
  N F_i\subseteq F_{i+1}\qquad\text{for }1\leq i<n.
\]
It sits between the ordinary Springer-fiber world and the broader Hessenberg
world, and it is important in the Schubert-calculus literature because of its
relationship with quantum cohomology of flag varieties.


\section[varietiesForKeys]{Demazure characters, keys, and flagged Schur polynomials}

Let $G=\operatorname{GL}_n(\setC)$ and let $B$ be the subgroup of upper
triangular matrices.  The flag variety can be written as the quotient $G/B$.
The subgroup $B$ is a \defin{Borel subgroup}\label{borelSubgroup}, and the irreducible closed
$B$-stable subvarieties inside $G/B$ are precisely Schubert varieties.  General
closed $B$-stable subvarieties are unions of Schubert varieties.

This is the geometric background for \hyperref[key]{key polynomials}.  A
dominant weight $\lambda$ gives a line bundle on $G/B$, and taking sections
over a Schubert variety gives a Demazure module.  Its character is a
\hyperref[demazureCharacter]{Demazure character}, also called a key polynomial
\cite{Demazure1974nouvelle,Demazure1974}.

The \hyperref[schurFlagged]{flagged Schur polynomials} are closely related.
They are Schur-type generating functions with row bounds, and they occur as
key polynomials in important cases.  In particular, every flagged Schur
polynomial is a key polynomial \cite[Thm.~23]{ReinerShimozono1995}, and
vexillary \hyperref[schubert]{Schubert polynomials} are flagged Schur
functions.


\section[varietiesFurtherReading]{Where to go next}

For computations in the cohomology of Grassmannians, see
\hyperref[schubertCalculus]{Schubert calculus}.  For the K-theoretic version,
see \hyperref[kTheory]{K-theory} and
\hyperref[grothendieck]{Grothendieck polynomials}.  For the matroid side of
Grassmannians, see \hyperref[positroid]{positroids} and
\hyperref[matroidBasePolytope]{matroid base polytopes}.  For Hessenberg
varieties and representation theory, see the \hyperref[dotAction]{dot action}
page.  For Springer-fiber and Richardson-tableau connections to insertion
algorithms, see the \hyperref[rsk]{RSK} page.
